\documentclass[12pt]{article}
\usepackage[margin=0.72in]{geometry}
\usepackage{lmodern}
\usepackage{microtype}
\usepackage{multicol}
\usepackage{fancyhdr}
\usepackage{titlesec}
\usepackage{enumitem}
\usepackage{xcolor}
\usepackage[hidelinks]{hyperref}

\setlength{\columnsep}{0.30in}
\setlength{\parindent}{1.05em}
\setlength{\parskip}{0.24em}
\setlength{\headheight}{15pt}
\titleformat{\section}{\large\bfseries\scshape}{}{0pt}{}
\titleformat{\subsection}{\normalsize\bfseries}{}{0pt}{}
\pagestyle{fancy}
\fancyhf{}
\fancyhead[L]{Aristotle on Responsibility}
\fancyhead[R]{Reading Handout}
\fancyfoot[C]{\thepage}
\renewcommand{\headrulewidth}{0.4pt}

\newcommand{\focusbox}[1]{%
\vspace{0.4em}\noindent\colorbox{gray!12}{\begin{minipage}{0.96\linewidth}#1\end{minipage}}\vspace{0.5em}}

\begin{document}
\begin{center}
{\LARGE\bfseries Are We Free, or Are We Ruled by Fate?}\par\vspace{0.25em}
{\large\scshape Reading 1: Aristotle on Voluntary Action}\par\vspace{0.45em}
\rule{0.82\textwidth}{0.4pt}\par\vspace{0.7em}
\end{center}

\section*{Vocabulary Preview}
\begin{multicols}{2}
\begin{itemize}[leftmargin=*, itemsep=0.18em]
\item \textbf{Voluntary}: Done knowingly, with the source of action in the person acting.
\item \textbf{Involuntary}: Done unwillingly because of force or ignorance, especially when the person later regrets it.
\item \textbf{Compulsion}: Being forced by an outside cause so that the person contributes nothing to the action.
\item \textbf{Ignorance}: Not knowing some important fact about what one is doing.
\item \textbf{Particular circumstances}: The concrete details of an action: who, what, to whom, with what, why, and how.
\item \textbf{Pardon}: Forgiveness or excuse because the action was not fully voluntary.
\item \textbf{Choice}: A deliberate desire for something within one's power.
\item \textbf{Deliberation}: Reasoning about what means will reach an end.
\item \textbf{Virtue}: A stable excellence of character formed by repeated action.
\item \textbf{Vice}: A bad habit or settled corruption of character.
\item \textbf{Character}: The kind of person one has become through habits and choices.
\item \textbf{Responsibility}: Being rightly open to praise, blame, reward, or punishment.
\end{itemize}
\end{multicols}

\section*{Introduction}
Aristotle was a Greek philosopher who lived from 384 to 322 B.C. In the \textit{Nicomachean Ethics}, he asks what kind of life is truly good for a human being. He argues that happiness is not merely a feeling, but a life of activity shaped by virtue, reason, and good habits.

Before Aristotle can explain virtue, he must answer a practical question: when should a person be praised or blamed? Not every harmful action deserves the same judgment. Some actions are forced. Some are done in ignorance. Some are chosen under pressure. Aristotle gives students a careful vocabulary for deciding when an action is truly ours.

\focusbox{\textbf{Source note:} Adapted and modernized from Aristotle, \textit{Nicomachean Ethics}, Book III, chapters 1--5, in F. H. Peters's public-domain translation. This version preserves Aristotle's sequence, examples, and main claims while using modern student-facing English.}

\begin{multicols}{2}
\section*{Reading}

\subsection*{1. Why Voluntary Action Matters}
Virtue has to do with feelings and actions. But praise and blame are given only for what is voluntary. What is involuntary receives pardon, and sometimes pity.

For that reason, anyone studying virtue must clearly distinguish voluntary actions from involuntary actions. Lawmakers also need this distinction, because rewards and punishments depend on whether a person is responsible.

In general, an action is involuntary when it is done because of \textbf{compulsion} or because of \textbf{ignorance}.

\subsection*{2. Actions Done by Force}
An action is compulsory when its cause is outside the person, and the person contributes nothing to it. For example, a person might be carried away by a whirlwind or dragged somewhere by people he cannot resist.

But some cases are harder. Suppose a tyrant orders you to do something shameful while holding your parents or children in his power. If you obey, they live. If you refuse, they die. Is your action voluntary or involuntary?

A similar case is throwing cargo overboard during a storm. No one would choose to lose his property for no reason. But a sensible person may throw it away to save himself and the crew.

Actions like these are mixed. In one sense, they are not what the person would choose by themselves. But at the actual time, under those circumstances, they are closer to voluntary actions. The source of the movement is still in the person. He moves his own body. He acts for a reason he sees at that moment. In such cases, it still rests with him to do or not do the action.

Therefore, these actions are voluntary at the time, though in themselves no one would choose them for their own sake.

\subsection*{3. Pressure Does Not Excuse Everything}
Sometimes people are praised for actions done under pressure. A person may endure something painful or shameful in order to secure a great and noble result. In other cases, people are blamed, especially if they endure disgrace for a small or bad reason.

Sometimes we do not praise, but pardon. A person may be pushed by pressure too strong for human nature to bear.

Still, Aristotle does not think the word \textit{compulsion} excuses everything. Some actions are so shameful that a person ought to suffer death rather than do them. It is often hard to decide whether we should do a certain deed to avoid a certain evil, or endure the evil rather than do the deed. It is even harder to stand by that decision when fear and pain arrive.

So, strictly speaking, a fully compulsory action is one whose cause lies outside the person, while the person contributes nothing to it. Mixed actions may deserve pity, praise, or blame depending on the circumstances and the reason for acting.

\subsection*{4. Actions Done Through Ignorance}
Actions done through ignorance are also not fully voluntary. But Aristotle makes an important distinction.

If a person acts through ignorance and is later pained and sorry when he discovers what he has done, the action is involuntary. If he is not sorry afterward, Aristotle calls the action \textit{not-voluntary}, but not fully involuntary.

Acting \textit{through ignorance} is different from acting \textit{in ignorance}. A drunk person or an angry person may not be thinking clearly. But Aristotle says the drunk man does not act through ignorance in the excusing sense. He acts through drunkenness. The angry man acts through anger. These conditions may cloud judgment, but they do not automatically remove responsibility.

A vicious person may also be ignorant of what should be done. But this kind of moral ignorance is part of his vice. It does not excuse him. Ignorance of what is right and wrong is blameworthy when the person should have learned better.

The ignorance that can make an action involuntary is ignorance of the particular facts: who is acting, what is being done, who is affected, what instrument is being used, why the action is being done, or how it is being done.

For example, someone might not know that a weapon is loaded. Someone might mistake a friend for an enemy. Someone might give a drug intending to heal, but the drug kills. In such cases, the person may not know the important particulars of the action.

The most important particulars are usually the person affected and the result. But even here, the person must be grieved and sorry afterward if the action is to be called involuntary.

\subsection*{5. Aristotle's Definition}
We can now define the voluntary.

A voluntary action is one that begins in the person acting, while the person knows the particular circumstances of the action.

This means anger and desire do not automatically make actions involuntary. If they did, animals and children would not act voluntarily at all. Also, we would absurdly claim that noble actions done from desire are voluntary, but shameful actions done from desire are involuntary, even though the cause is the same.

It is also strange to blame outside things for our disgraceful actions while claiming our noble actions as our own. The things we desire may attract us, but our readiness to yield to them belongs to us.

For Aristotle, the passions belong to the person just as reason does. Actions done from anger or desire are still the person's actions.

\subsection*{6. Choice and Deliberation}
After distinguishing voluntary from involuntary action, Aristotle turns to choice. Choice is closely connected with virtue and is an even clearer test of character than action alone.

Choice is not exactly the same as desire, anger, wish, or opinion. We can desire things without choosing them. We can wish for impossible things, such as never dying, but we cannot choose the impossible. We choose only what we think can be brought about by our own action.

Choice involves deliberation. We do not deliberate about everything. No one deliberates about eternal things, such as the stars, or things that happen by necessity, such as sunrise. We deliberate about conduct that is within our power.

We also deliberate about means, not ends. A doctor does not deliberate about whether to heal; healing is the end of medicine. The doctor deliberates about how to heal. In the same way, a person who wants a good end considers the means by which it can be reached.

Aristotle defines choice as deliberate desire for something in our power. We deliberate first, then decide, then desire the chosen action in light of that decision.

\subsection*{7. Are Virtue and Vice Up to Us?}
Actions guided by choice are voluntary. Since virtues show themselves in actions, virtue depends on us. Vice also depends on us.

Where it lies with us to do something, it also lies with us not to do it. Where we can say no, we can say yes. If doing a noble deed lies within our power, then failing to do it also lies within our power. If avoiding a shameful deed lies within our power, then doing that shameful deed also lies within our power.

Therefore, Aristotle argues, it lies with us to become worthy or worthless people.

This does not mean change is easy. A person forms character by repeated actions. People become unjust by doing unjust acts. They become self-indulgent by repeatedly giving themselves over to pleasure. At first, the beginning is in their power.

Aristotle compares this to throwing a stone. Once you throw it, you cannot call it back. But the throwing was in your power at the beginning. In the same way, a person who has built a bad character may not be able to change instantly whenever he wishes. But he was responsible for the repeated choices that formed that character.

So Aristotle's view is demanding: our actions form our habits, our habits form our character, and our character shapes what seems good to us. Because of this, responsibility does not disappear simply because desire, anger, habit, or pressure influence us.

\section*{Reading Focus}
\begin{enumerate}[leftmargin=*]
\item In one sentence, define Aristotle's rule for when an action is voluntary.
\item Why does Aristotle think mixed actions, such as throwing cargo overboard in a storm, are closer to voluntary than involuntary?
\item What is the difference between acting \textit{through ignorance} and acting \textit{in ignorance}?
\item Aristotle says anger and desire do not automatically excuse us. Do you agree? Give one example.
\item Explain the stone example. What does it suggest about habit and responsibility?
\end{enumerate}

\section*{Window Note and Summary}
Create a window note that visually shows Aristotle's three categories: voluntary action, involuntary action, and mixed action. Include at least one drawing or symbol for each category. At the bottom, write a 3--5 sentence summary explaining how Aristotle would decide whether a person deserves praise, blame, pardon, or pity.

\section*{Connection to the Unit Question}
Write one claim answering this question: \textbf{If our actions are shaped by fear, desire, ignorance, and habit, are they still ours?} Use Aristotle to support your answer.
\end{multicols}
\end{document}
